% !TEX root = ms.tex

\documentclass[a4paper,12pt, oneside]{book}
\usepackage[utf8]{inputenc}
\usepackage{xcolor}
\usepackage{titlesec}
\usepackage{graphicx}
\usepackage{hyperref}
\usepackage{amsmath}
\usepackage{natbib}
\usepackage{lipsum}
\usepackage{titling}
\usepackage{graphicx}
\usepackage{titlesec}
\usepackage{fancyhdr}
\usepackage[nottoc]{tocbibind}
\usepackage{mathptmx}
\usepackage{transparent}
\usepackage{soul}
\usepackage{amssymb}
\usepackage{booktabs}
\usepackage{array}
\usepackage{caption}
\usepackage{pgfplots}
\usepackage{pgfplotstable}
\usepackage{adjustbox}
\usepackage[T1]{fontenc}
\usepackage[utf8]{inputenc}
\usepackage[spanish]{babel}
\usepackage{float}
\usepackage{newfloat}
\usepackage{setspace}
\usepackage{listings}


\titleformat{\chapter}[hang]
    {\fontsize{30}{40}\scshape\sffamily}
    {\thechapter .}
    {15pt}
    {}
    [\titlerule]

\titlespacing*{\chapter}{0pt}{0pt}{10pt}

\definecolor{lightblue}{rgb}{0.85, 0.96, 0.98}
\definecolor{lightred}{rgb}{0.99, 0.91, 0.92}
\definecolor{lightpurple}{rgb}{0.87, 0.72, 0.85}

\newcommand{\hlblue}[1]{{\sethlcolor{lightblue}\hl{#1}}}
\newcommand{\hlred}[1]{{\sethlcolor{lightred}\hl{#1}}}
\newcommand{\hlpurple}[1]{{\sethlcolor{lightpurple}\hl{#1}}}

\newcommand{\fix}[2]{
  \textcolor{red}{\hlred{[#1]}}
  \textcolor{blue}{\hlblue{(#2)}}
}

\newcommand{\todo}[1]{
	\textcolor{violet}{\transparent{0.7}\texttt{\#TODO:} \textsf{#1}}
}

\newcommand{\note}[2]{
	\textcolor{blue}{\transparent{0.7}\texttt{Note:} \textsf{#1}}
}

\pgfplotsset{compat=1.18}
\setlength{\parskip}{1em}
\setlength{\parindent}{0.5cm}
\singlespacing

\pagestyle{fancy}
\fancyhf{}
\fancyfoot[C]{\thepage}
\fancyhead[R]{\transparent{0.5}{\nouppercase{{\rightmark}}}}
\renewcommand{\headrulewidth}{0pt}
\renewcommand{\chaptermark}[1]{\markboth{#1}{}}

\renewcommand{\arraystretch}{1.3}
\setlength{\tabcolsep}{8pt}

\usepackage[utf8]{inputenc}   % sólo una vez
\usepackage[T1]{fontenc}
\usepackage[spanish]{babel}

\usepackage{float}
\usepackage{newfloat}
\DeclareFloatingEnvironment[
  fileext=lol,
  listname=Lista de códigos,
  name=Código
]{code}

\usepackage{caption}
\captionsetup[code]{font=small, labelfont=bf}
\usepackage{xcolor}
\usepackage{listings}
\lstset{
  language=Python,
  basicstyle={\fontfamily{lmvtt}\selectfont\footnotesize\hspace{-0.2em}},
  keywordstyle=\color[rgb]{0.86, 0.20, 0.44},       % rosa fuerte (keywords)
  stringstyle=\color[rgb]{0.60, 0.44, 0.12},        % naranja (strings)
  commentstyle=\color[rgb]{0.50, 0.50, 0.50},       % gris (comentarios)
  numbers=left,
  numberstyle=\tiny\color{gray},
  stepnumber=1,
  numbersep=10pt,
  frame=single,
  breaklines=true,
  showstringspaces=false,
  tabsize=4,
  inputencoding=utf8,
  extendedchars=true
}

